\nonstopmode{}
\documentclass[a4paper]{book}
\usepackage[times,inconsolata,hyper]{Rd}
\usepackage{makeidx}
\usepackage[utf8]{inputenc} % @SET ENCODING@
% \usepackage{graphicx} % @USE GRAPHICX@
\makeindex{}
\begin{document}
\chapter*{}
\begin{center}
{\textbf{\huge \R{} documentation}} \par\bigskip{{\Large of all in \file{man}}}
\par\bigskip{\large \today}
\end{center}
\Rdcontents{\R{} topics documented:}
\inputencoding{utf8}
\HeaderA{audit\_model\_registry}{Audit every registered scale and model prompt.}{audit.Rul.model.Rul.registry}
%
\begin{Description}
Audit every registered scale and model prompt.
\end{Description}
%
\begin{Usage}
\begin{verbatim}
audit_model_registry(
  registry_dir = NULL,
  output_dir = NULL,
  apply_unifications = FALSE
)
\end{verbatim}
\end{Usage}
%
\begin{Arguments}
\begin{ldescription}
\item[\code{registry\_dir}] Optional registry root. Defaults to bundled prompts.

\item[\code{output\_dir}] Optional directory for audit and unification CSV reports.

\item[\code{apply\_unifications}] Move exact duplicate prompt folders to a timestamped
backup directory. Defaults to \code{FALSE}.
\end{ldescription}
\end{Arguments}
%
\begin{Value}
A list with checks, duplicate\_groups, unification\_plan, and backup\_dir.
\end{Value}
\inputencoding{utf8}
\HeaderA{available\_provider\_models}{List models currently available to the configured provider.}{available.Rul.provider.Rul.models}
%
\begin{Description}
List models currently available to the configured provider.
\end{Description}
%
\begin{Usage}
\begin{verbatim}
available_provider_models(
  provider,
  api_key = NULL,
  timeout = 30,
  supported_only = TRUE
)
\end{verbatim}
\end{Usage}
%
\begin{Arguments}
\begin{ldescription}
\item[\code{provider}] \code{"gemini"}, \code{"openai"}, \code{"claude"}, or an alias.

\item[\code{api\_key}] Optional in-memory API key. Otherwise the provider's usual
environment variables are used.

\item[\code{timeout}] Request timeout in seconds.

\item[\code{supported\_only}] For Gemini, keep models supporting \code{generateContent}.
\end{ldescription}
\end{Arguments}
%
\begin{Value}
A data frame with provider, model id, display name, and raw metadata.
\end{Value}
\inputencoding{utf8}
\HeaderA{available\_scales}{List scales available in the local prompt registry.}{available.Rul.scales}
%
\begin{Description}
List scales available in the local prompt registry.
\end{Description}
%
\begin{Usage}
\begin{verbatim}
available_scales(registry_dir = NULL)
\end{verbatim}
\end{Usage}
%
\begin{Arguments}
\begin{ldescription}
\item[\code{registry\_dir}] Optional registry root. Defaults to bundled package prompts.
\end{ldescription}
\end{Arguments}
\inputencoding{utf8}
\HeaderA{build\_scale\_prompt}{Append article text to a scale prompt.}{build.Rul.scale.Rul.prompt}
%
\begin{Description}
Append article text to a scale prompt.
\end{Description}
%
\begin{Usage}
\begin{verbatim}
build_scale_prompt(prompt_text, article_text)
\end{verbatim}
\end{Usage}
%
\begin{Arguments}
\begin{ldescription}
\item[\code{prompt\_text}] Scale prompt text.

\item[\code{article\_text}] Extracted article text.
\end{ldescription}
\end{Arguments}
\inputencoding{utf8}
\HeaderA{compare\_training\_iterations}{Compare local prompt-training iterations with a reviewed answer key.}{compare.Rul.training.Rul.iterations}
%
\begin{Description}
Each iteration directory must contain one timestamped \AsIs{\texttt{*\_Consensus\_Report.csv}}
produced by \code{\LinkA{run\_dataset()}{run.Rul.dataset}}. The answer key must contain an \code{ID} column and
columns named \code{IT01}, \code{IT02}, etc., or another \code{IT}-prefixed item scheme.
\end{Description}
%
\begin{Usage}
\begin{verbatim}
compare_training_iterations(ideal_csv, iterations_dir)
\end{verbatim}
\end{Usage}
%
\begin{Arguments}
\begin{ldescription}
\item[\code{ideal\_csv}] Path to the reviewed answer-key CSV.

\item[\code{iterations\_dir}] Directory containing \code{iteration\_1}, \code{iteration\_2}, etc.
\end{ldescription}
\end{Arguments}
%
\begin{Value}
A list with \code{comparison} (item-level rows) and \code{summary} (accuracy by
iteration).
\end{Value}
\inputencoding{utf8}
\HeaderA{convert\_article\_markdown\_llm}{Convert extracted article text to faithful Markdown with an LLM.}{convert.Rul.article.Rul.markdown.Rul.llm}
%
\begin{Description}
Convert extracted article text to faithful Markdown with an LLM.
\end{Description}
%
\begin{Usage}
\begin{verbatim}
convert_article_markdown_llm(
  article_text,
  provider = "openai",
  model = "gpt-5.6-luna",
  temperature = 0,
  prompt = NULL,
  max_chars = 50000,
  ...
)
\end{verbatim}
\end{Usage}
%
\begin{Arguments}
\begin{ldescription}
\item[\code{article\_text}] Extracted article text.

\item[\code{provider}] LLM provider.

\item[\code{model}] Provider model.

\item[\code{temperature}] Sampling temperature.

\item[\code{prompt}] Optional local conversion prompt.

\item[\code{...}] Additional provider call settings.
\end{ldescription}
\end{Arguments}
\inputencoding{utf8}
\HeaderA{extract\_article\_text}{Extract text from a supported article file.}{extract.Rul.article.Rul.text}
%
\begin{Description}
Extract text from a supported article file.
\end{Description}
%
\begin{Usage}
\begin{verbatim}
extract_article_text(
  file_path,
  filetype = "auto",
  strip_references = TRUE,
  tables_advanced = TRUE,
  cache_markdown = TRUE,
  conversion = "basic",
  provider = "gemini",
  model = "gemini-3.6-flash",
  conversion_prompt = NULL,
  temperature = 0,
  max_chars = 50000,
  ...
)
\end{verbatim}
\end{Usage}
%
\begin{Arguments}
\begin{ldescription}
\item[\code{file\_path}] Path to a \code{.pdf}, \code{.txt}, or \code{.md} article file.

\item[\code{filetype}] One of \code{"auto"}, \code{"pdf"}, \code{"txt"}, or \code{"md"}.

\item[\code{strip\_references}] Whether to remove the reference section before sending text to an LLM.

\item[\code{conversion}] PDF conversion mode passed to \code{extract\_pdf\_text()}.

\item[\code{provider, model, conversion\_prompt, temperature}] LLM conversion settings.

\item[\code{max\_chars}] Maximum text size per LLM conversion call.
\end{ldescription}
\end{Arguments}
%
\begin{Details}
Files are read from the local filesystem. PDF, TXT, and Markdown
inputs are supported.
\end{Details}
\inputencoding{utf8}
\HeaderA{extract\_pdf\_text}{Extract text from a PDF article.}{extract.Rul.pdf.Rul.text}
%
\begin{Description}
Extract text from a PDF article.
\end{Description}
%
\begin{Usage}
\begin{verbatim}
extract_pdf_text(
  pdf_path,
  strip_references = TRUE,
  tables_advanced = TRUE,
  cache_markdown = TRUE,
  conversion = "basic",
  provider = "openai",
  model = "gpt-5.6-luna",
  conversion_prompt = NULL,
  temperature = 0,
  max_chars = 50000,
  ...
)
\end{verbatim}
\end{Usage}
%
\begin{Arguments}
\begin{ldescription}
\item[\code{pdf\_path}] Path to a PDF article.

\item[\code{strip\_references}] Whether to remove the reference section before sending text to an LLM.

\item[\code{tables\_advanced}] Whether to apply local table heuristics.

\item[\code{cache\_markdown}] Whether to cache PDF conversion beside the source PDF.

\item[\code{conversion}] PDF conversion mode: \code{"basic"} or \code{"llm"}. LLM conversion
usually produces better reading order and Markdown structure for multi-column articles.

\item[\code{provider, model, conversion\_prompt, temperature}] LLM conversion settings.

\item[\code{max\_chars}] Maximum extracted-text size per LLM conversion call. Larger
PDFs are split into line-safe chunks and reassembled into one Markdown file.
\end{ldescription}
\end{Arguments}
\inputencoding{utf8}
\HeaderA{parse\_scale\_scores}{Parse scale item scores from an LLM response.}{parse.Rul.scale.Rul.scores}
%
\begin{Description}
Parse scale item scores from an LLM response.
\end{Description}
%
\begin{Usage}
\begin{verbatim}
parse_scale_scores(text, items = 1:10, metadata = NULL)
\end{verbatim}
\end{Usage}
%
\begin{Arguments}
\begin{ldescription}
\item[\code{text}] LLM response text.

\item[\code{items}] Numeric item ids to parse. Defaults to 1:10.

\item[\code{metadata}] Optional resolved scale metadata.
\end{ldescription}
\end{Arguments}
\inputencoding{utf8}
\HeaderA{propose\_prompt\_revision}{Propose a revised prompt from the results of a training iteration.}{propose.Rul.prompt.Rul.revision}
%
\begin{Description}
This function does not modify files, execute articles, or replace the
current prompt. It asks the selected provider to propose a complete revised
prompt using the current prompt, item-level comparison, and audit logs from
the latest iteration.
\end{Description}
%
\begin{Usage}
\begin{verbatim}
propose_prompt_revision(
  prompt_path,
  comparison,
  reason_files,
  provider = "gemini",
  model = "gemini-3.6-flash",
  api_key = NULL,
  project_id = NULL,
  temperature = 0,
  timeout = 300,
  output_path = NULL
)
\end{verbatim}
\end{Usage}
%
\begin{Arguments}
\begin{ldescription}
\item[\code{prompt\_path}] Path to the current prompt.md.

\item[\code{comparison}] A comparison data frame, or path to comparison.csv.

\item[\code{reason\_files}] Character vector of audit-log or evidence files.

\item[\code{provider}] LLM provider used to generate the proposal.

\item[\code{model}] Model used to generate the proposal.

\item[\code{api\_key}] Optional in-memory API key.

\item[\code{project\_id}] Optional OpenAI project identifier.

\item[\code{temperature}] Generation temperature. Defaults to 0.

\item[\code{timeout}] Request timeout in seconds.

\item[\code{output\_path}] Optional path where the proposal is written as Markdown.
\end{ldescription}
\end{Arguments}
%
\begin{Value}
A list with the proposed prompt, source paths, provider, and model.
\end{Value}
\inputencoding{utf8}
\HeaderA{resolve\_prompt}{Resolve the best prompt registered for a scale and requested model.}{resolve.Rul.prompt}
%
\begin{Description}
Resolution order is exact model, same model family token, same numeric version,
The accepted prompt is always the scale-level prompt in
\code{scales/<scale>/prompt.md}, with its metadata in the sibling
\code{metadata.json}. No model- or provider-specific fallback is used.
\end{Description}
%
\begin{Usage}
\begin{verbatim}
resolve_prompt(scale, model, provider = NULL, registry_dir = NULL)
\end{verbatim}
\end{Usage}
%
\begin{Arguments}
\begin{ldescription}
\item[\code{scale}] Scale name, for example \code{"mqs"}.

\item[\code{model}] Requested model name.

\item[\code{provider}] Optional provider used only for the API call. It does not
select a prompt.

\item[\code{registry\_dir}] Optional registry root. Defaults to bundled package prompts.
\end{ldescription}
\end{Arguments}
%
\begin{Value}
A list containing prompt path, model folder, match strategy, and metadata.
\end{Value}
\inputencoding{utf8}
\HeaderA{run\_article}{Run one article file through a registered scale prompt.}{run.Rul.article}
%
\begin{Description}
Run one article file through a registered scale prompt.
\end{Description}
%
\begin{Usage}
\begin{verbatim}
run_article(
  article_path = NULL,
  scale = "mqs",
  provider = "gemini",
  model = "gemini-3.6-flash",
  registry_dir = NULL,
  filetype = "auto",
  strip_references = TRUE,
  tables_advanced = TRUE,
  cache_markdown = TRUE,
  conversion = "basic",
  conversion_provider = "openai",
  conversion_model = "gpt-5.6-luna",
  conversion_prompt = NULL,
  max_chars = 50000,
  items = NULL,
  output_dir = NULL,
  temperature = 0,
  top_p = 0.1,
  timeout = 300,
  api_key = NULL,
  project_id = NULL,
  pdf_path = NULL,
  reasoning_effort = NULL,
  max_retries = 3,
  retry_wait_seconds = 1,
  retry_backoff = 2,
  rate_limit_seconds = 0,
  reference_scores = NULL,
  validation_mode = c("free", "reference"),
  write_evidence = TRUE
)
\end{verbatim}
\end{Usage}
%
\begin{Arguments}
\begin{ldescription}
\item[\code{article\_path}] Article path. Supported extensions are \code{.pdf}, \code{.txt}, and \code{.md}.

\item[\code{scale}] Scale name. Defaults to \code{"mqs"}.

\item[\code{provider}] LLM provider. Defaults to \code{"gemini"}.

\item[\code{model}] Requested model. Defaults to \code{"gemini-3.6-flash"}.

\item[\code{registry\_dir}] Optional prompt registry root.

\item[\code{filetype}] One of \code{"auto"}, \code{"pdf"}, \code{"txt"}, or \code{"md"}.

\item[\code{strip\_references}] Whether to remove references before the LLM call.

\item[\code{tables\_advanced}] Whether to structure extracted headings, lists, and
table-like blocks as Markdown before the LLM call.

\item[\code{cache\_markdown}] Whether PDF extraction writes a same-name \code{.md} cache
beside the source PDF for reuse by \code{filetype = "auto"}.

\item[\code{conversion}] PDF conversion mode: \code{"basic"} or \code{"llm"}. LLM conversion
generally gives better results for multi-column articles and complex layouts.

\item[\code{conversion\_provider, conversion\_model, conversion\_prompt}] Settings for the
LLM used only to convert PDF text into Markdown.

\item[\code{max\_chars}] Maximum text size per LLM PDF-conversion call; larger PDFs
are split and reassembled automatically.

\item[\code{items}] Item ids to parse. Defaults to 1:10.

\item[\code{output\_dir}] Optional output directory for an audit log.

\item[\code{temperature, top\_p, timeout, api\_key, project\_id}] Provider call settings.

\item[\code{pdf\_path}] Backward-compatible alias for \code{article\_path}.

\item[\code{max\_retries, retry\_wait\_seconds, retry\_backoff, rate\_limit\_seconds}] Reliability settings.

\item[\code{reference\_scores}] Optional reviewed scores for reference validation.

\item[\code{validation\_mode}] Either \code{"free"} or \code{"reference"}.

\item[\code{write\_evidence}] Whether to write the per-article evidence CSV. Dataset
workflows set this to \code{FALSE} because they write one consolidated report.
\end{ldescription}
\end{Arguments}
\inputencoding{utf8}
\HeaderA{run\_dataset}{Run article files in a directory through a registered scale prompt.}{run.Rul.dataset}
%
\begin{Description}
Existing \AsIs{\texttt{*\_AuditLog.txt}} files larger than 100 bytes are reused and skipped.
\end{Description}
%
\begin{Usage}
\begin{verbatim}
run_dataset(
  articles_dir,
  scale = "mqs",
  provider = "gemini",
  model = "gemini-3.6-flash",
  output_dir,
  registry_dir = NULL,
  filetype = "auto",
  strip_references = TRUE,
  tables_advanced = TRUE,
  cache_markdown = TRUE,
  conversion = "basic",
  conversion_provider = "openai",
  conversion_model = "gpt-5.6-luna",
  conversion_prompt = NULL,
  max_articles = 0,
  max_chars = 50000,
  items = NULL,
  temperature = 0,
  top_p = 0.1,
  timeout = 300,
  api_key = NULL,
  project_id = NULL,
  max_retries = 3,
  retry_wait_seconds = 1,
  retry_backoff = 2,
  rate_limit_seconds = 0,
  reference_csv = NULL,
  reasoning_effort = NULL,
  validation_mode = c("free", "reference")
)
\end{verbatim}
\end{Usage}
%
\begin{Arguments}
\begin{ldescription}
\item[\code{articles\_dir}] Directory containing articles.

\item[\code{scale}] Scale name.

\item[\code{provider}] LLM provider.

\item[\code{model}] Requested model.

\item[\code{output\_dir}] Directory where audit logs and CSV are written.

\item[\code{registry\_dir}] Optional prompt registry root.

\item[\code{filetype}] One of \code{"auto"}, \code{"pdf"}, \code{"txt"}, or \code{"md"}.

\item[\code{strip\_references, items, temperature, top\_p, timeout, api\_key, project\_id}] Provider and parsing settings.

\item[\code{conversion, conversion\_provider, conversion\_model, conversion\_prompt}] PDF
conversion settings. LLM conversion generally gives better results for
multi-column articles and complex layouts.

\item[\code{max\_articles}] Maximum pending PDFs to process. \code{0} means all.

\item[\code{max\_chars}] Maximum text size per LLM PDF-conversion call.

\item[\code{max\_retries, retry\_wait\_seconds, retry\_backoff, rate\_limit\_seconds}] Reliability settings.

\item[\code{reference\_csv}] Optional reviewed-score CSV.

\item[\code{validation\_mode}] Either \code{"free"} or \code{"reference"}.
\end{ldescription}
\end{Arguments}
%
\begin{Value}
A list containing the timestamped run directory, consensus report,
consolidated evidence report, and errors when any occurred.
\end{Value}
\inputencoding{utf8}
\HeaderA{run\_llm}{Call a supported LLM provider with user-owned API credentials.}{run.Rul.llm}
%
\begin{Description}
Call a supported LLM provider with user-owned API credentials.
\end{Description}
%
\begin{Usage}
\begin{verbatim}
run_llm(
  prompt,
  provider = "gemini",
  model = "gemini-3.6-flash",
  temperature = 0,
  top_p = 0.1,
  timeout = 300,
  api_key = NULL,
  project_id = NULL,
  max_retries = 3,
  retry_wait_seconds = 1,
  retry_backoff = 2,
  rate_limit_seconds = 0,
  response_schema = NULL,
  reasoning_effort = NULL
)
\end{verbatim}
\end{Usage}
%
\begin{Arguments}
\begin{ldescription}
\item[\code{prompt}] Prompt text to send.

\item[\code{provider}] \code{"gemini"}, \code{"openai"}, \code{"chatgpt"}, \code{"claude"}, or \code{"anthropic"}.

\item[\code{model}] Model id.

\item[\code{temperature}] Sampling temperature. For GPT-5.6 models it is sent only
when \code{reasoning\_effort = "none"}.

\item[\code{top\_p}] Gemini nucleus-sampling value; ignored by OpenAI and Claude.

\item[\code{timeout}] Maximum duration in seconds for each individual API attempt.

\item[\code{api\_key}] Optional in-memory API key.

\item[\code{project\_id}] Optional OpenAI project id.

\item[\code{max\_retries}] Maximum retries for transient failures.

\item[\code{retry\_wait\_seconds}] Initial exponential-backoff delay.

\item[\code{retry\_backoff}] Multiplicative exponential-backoff factor.

\item[\code{rate\_limit\_seconds}] Minimum delay between requests in this R process.

\item[\code{response\_schema}] Optional registered response schema.

\item[\code{reasoning\_effort}] Optional Responses API reasoning effort, for example
\code{"none"}, \code{"low"}, or \code{"medium"}. GPT-5.6 models require
\code{"none"} when \code{temperature} is used; otherwise temperature is omitted.
\end{ldescription}
\end{Arguments}
\inputencoding{utf8}
\HeaderA{strip\_references\_section}{Remove reference sections from extracted article text.}{strip.Rul.references.Rul.section}
%
\begin{Description}
Remove reference sections from extracted article text.
\end{Description}
%
\begin{Usage}
\begin{verbatim}
strip_references_section(article_text)
\end{verbatim}
\end{Usage}
%
\begin{Arguments}
\begin{ldescription}
\item[\code{article\_text}] Full extracted text.
\end{ldescription}
\end{Arguments}
\inputencoding{utf8}
\HeaderA{structure\_article\_markdown}{Convert extracted text to Markdown headings, lists, and tables.}{structure.Rul.article.Rul.markdown}
%
\begin{Description}
Convert extracted text to Markdown headings, lists, and tables.
\end{Description}
%
\begin{Usage}
\begin{verbatim}
structure_article_markdown(article_text, tables_advanced = TRUE)
\end{verbatim}
\end{Usage}
%
\begin{Arguments}
\begin{ldescription}
\item[\code{article\_text}] Extracted article text.

\item[\code{tables\_advanced}] Whether to convert table-like blocks.
\end{ldescription}
\end{Arguments}
\printindex{}
\end{document}
